\documentclass[11pt,a4paper]{article}

\usepackage[T1]{fontenc}
\usepackage[utf8]{inputenc}
\usepackage[italian]{babel}
\usepackage[a4paper,margin=2.2cm]{geometry}
\usepackage{hyperref}
\usepackage{booktabs}
\usepackage{longtable}
\usepackage{tabularx}
\usepackage{array}
\usepackage{enumitem}
\usepackage{xcolor}

\hypersetup{
  colorlinks=true,
  linkcolor=blue!50!black,
  urlcolor=blue!50!black,
  pdftitle={HP Fortress - Technical Documentation},
  pdfauthor={Alessandro Mainardi}
}

\setlist[itemize]{leftmargin=1.4em}
\setlist[enumerate]{leftmargin=1.8em}

\title{HP Fortress\\\large Technical Documentation}
\author{Alessandro Mainardi}
\date{Versione 2.0 \textbar{} \today}

\begin{document}

\maketitle
\tableofcontents
\newpage

\section{Project Overview}

\subsection{Scopo del Progetto}
Il progetto \textbf{HP Fortress} realizza un'infrastruttura server on-premise ibrida per due obiettivi paralleli:
\begin{enumerate}
  \item \textbf{Educational \& Engineering Lab}: ambiente pratico per Windows Server 2025, Active Directory, Hyper-V, Linux/Docker e SQL Server.
  \item \textbf{Home Production Environment}: NAS affidabile per backup familiari e servizi multimediali self-hosted.
\end{enumerate}

L'architettura è progettata per funzionamento 24/7, con segregazione dei ruoli tra host, VM infrastrutturali e servizi applicativi.

\subsection{Requisiti Funzionali}
\begin{itemize}
  \item \textbf{Storage resiliente}: mirror software (RAID 1) con ReFS.
  \item \textbf{Virtualizzazione}: VM Windows/Linux in esecuzione concorrente.
  \item \textbf{Backup centralizzato}: condivisioni SMB con ACL granulari.
  \item \textbf{Media streaming}: Plex con supporto a transcodifica hardware.
  \item \textbf{Lab di sviluppo}: SQL Server e IIS per workload di test/portfolio.
  \item \textbf{Accesso remoto sicuro}: VPN/Tunnel senza esposizione diretta di porte critiche.
\end{itemize}

\subsection{Vincoli e Presupposti}
\begin{itemize}
  \item Budget hardware target inferiore a EUR 600.
  \item Base hardware: workstation enterprise ricondizionata (HP EliteDesk 800 G4 Tower).
  \item Licensing software: Azure Dev Tools for Teaching (uso didattico/non commerciale diretto).
  \item Ambiente domestico: attenzione a rumore, consumi e gestione termica.
\end{itemize}

\paragraph{Rischio accettato (power)}
Assenza di UPS: blackout o brownout possono introdurre rischio di corruzione dati. Mitigazione: backup frequenti, ripristino testato e politiche conservative di riavvio/alimentazione.

\section{Hardware Architecture (Physical Layer)}

\subsection{Bill of Materials (BOM)}
\begin{longtable}{>{\raggedright\arraybackslash}p{0.22\textwidth} >{\raggedright\arraybackslash}p{0.44\textwidth} >{\raggedright\arraybackslash}p{0.14\textwidth} >{\raggedright\arraybackslash}p{0.12\textwidth}}
\toprule
\textbf{Componente} & \textbf{Modello / Specifica} & \textbf{Ruolo} & \textbf{Stato} \\
\midrule
\endfirsthead
\toprule
\textbf{Componente} & \textbf{Modello / Specifica} & \textbf{Ruolo} & \textbf{Stato} \\
\midrule
\endhead
Compute Node & HP EliteDesk 800 G4 Tower (Q370) & Host Hyper-V & Refurbished \\
CPU & Intel Core \textbf{i7-8700K} (6C/12T) & Compute & Included \\
RAM & 16 GB DDR4 2666 MHz (2x8 GB) & Base memory & Included \\
Boot Storage & 2x NVMe SSD 512 GB & OS host/VM system disks & Included \\
Data Storage 1 & WD Red Plus 4 TB (WD40EFPX, CMR) & Mirror A & New \\
Data Storage 2 & WD Red Plus 4 TB (WD40EFPX, CMR) & Mirror B & New \\
Accessori & Cavi SATA, viti/grommet, pasta termica & Assembly & New \\
\bottomrule
\end{longtable}

\subsection{Specifiche Tecniche del Nodo}
\begin{itemize}
  \item CPU: Intel Core i7-8700K.
  \item RAM installata: 16 GB DDR4 dual-channel (espandibile).
  \item Storage raw: circa 9 TB (2x512 GB NVMe + 2x4 TB HDD).
  \item NIC onboard: Intel I219-LM 1GbE.
  \item Espansioni: slot PCIe x16/x4, PCIe x1 e 2 slot M.2 2280.
\end{itemize}

\subsection{BIOS \& Firmware Configuration (UEFI)}
\begin{itemize}
  \item Secure Boot: enabled; legacy support: disabled.
  \item Intel VT-x e VT-d: enabled.
  \item Modalità storage: AHCI (nessun RAID BIOS).
  \item Power profile: configurazione orientata a stabilità/prestazioni.
  \item Power loss policy: coerente con strategia operativa locale.
\end{itemize}

\subsection{Port Map e Storage Topology}
\begin{itemize}
  \item SATA 1: WD Red Plus 4 TB (mirror member A).
  \item SATA 2: WD Red Plus 4 TB (mirror member B).
  \item M.2 slot 1: NVMe SSD 512 GB (boot/host).
  \item M.2 slot 2: NVMe SSD 512 GB (boot/VM support).
\end{itemize}

\section{Host System Configuration (Hypervisor Layer)}

\subsection{Operating System}
\begin{itemize}
  \item OS host: Windows Server 2025 Datacenter (Desktop Experience).
  \item Ruolo: Hyper-V host + file services/storage.
  \item Hostname: \texttt{HV-NODE-01}.
  \item Principio operativo: host minimale, servizi applicativi nelle VM.
\end{itemize}

\subsection{Network Configuration (Host)}
\begin{itemize}
  \item NIC fisica: Intel I219-LM con driver compatibile.
  \item Strategia IP: statica.
  \item vSwitch Hyper-V: \texttt{vSwitch\_External} in modalità External.
  \item DNS host: priorità DNS interno (DC) con fallback pubblico controllato.
\end{itemize}

\subsection{Storage Strategy (Software Defined Storage)}
\begin{itemize}
  \item Storage Pool: \texttt{SP\_DATA\_01} su 2 HDD WD Red Plus.
  \item Virtual disk: \texttt{vDisk\_Data}, layout mirror a 2 vie.
  \item Volume dati: lettera \texttt{D:}, file system ReFS.
  \item Feature: integrity streams; dedup opzionale su dataset idonei.
\end{itemize}

\subsection{Security Hardening (Host Level)}
\begin{itemize}
  \item Firewall host attivo con regole inbound minime.
  \item RDP consentito solo ad amministratori.
  \item Aggiornamenti pianificati e riavvio controllato.
  \item BitLocker su volumi sensibili con gestione sicura recovery key.
\end{itemize}

\section{File System \& Data Governance}

\subsection{Directory Structure (Volume D:)}
Struttura logica principale:
\begin{itemize}
  \item \texttt{D:\textbackslash{}MEDIA} (Movies, TV\_Shows, Music)
  \item \texttt{D:\textbackslash{}FAMILY\_BACKUP} (cartelle private per utente)
  \item \texttt{D:\textbackslash{}DEV} (ISO, VMs, Repos)
  \item \texttt{D:\textbackslash{}PUBLIC} (scambio temporaneo)
  \item \texttt{D:\textbackslash{}APP\_DATA} (persistenza applicazioni container)
\end{itemize}

\subsection{Identity \& Access Management (IAM)}
\begin{itemize}
  \item Utenti consumer (famiglia) definiti come utenti locali per resilienza operativa.
  \item Service account dedicati per Plex/Immich.
  \item ACL NTFS per privacy per-cartella su \texttt{FAMILY\_BACKUP}.
  \item Ereditarietà disabilitata dove necessario per isolamento rigoroso.
\end{itemize}

\subsection{Sharing Protocols}
\begin{itemize}
  \item Share SMB principali: \texttt{MEDIA}, \texttt{FAMILY\_BACKUP}.
  \item Access-Based Enumeration (ABE) abilitata.
  \item Network discovery e profilo private configurati per rete locale.
\end{itemize}

\section{Virtualization Strategy (Workloads)}

\subsection{VM Inventory}
\begin{longtable}{>{\raggedright\arraybackslash}p{0.24\textwidth} >{\raggedright\arraybackslash}p{0.20\textwidth} >{\raggedright\arraybackslash}p{0.30\textwidth} >{\raggedright\arraybackslash}p{0.14\textwidth}}
\toprule
\textbf{VM} & \textbf{OS} & \textbf{Ruolo} & \textbf{Priorità} \\
\midrule
\endfirsthead
\toprule
\textbf{VM} & \textbf{OS} & \textbf{Ruolo} & \textbf{Priorità} \\
\midrule
\endhead
\texttt{VM-DC-01} & Windows Server 2025 & Domain Controller, DNS, DHCP & Alta \\
\texttt{VM-LINUX-01} & Ubuntu Server 24.04 LTS & Docker host servizi media/backup/tunnel & Media \\
\texttt{VM-SQL-01} & Windows Server 2025 & SQL Server e IIS (lab/portfolio) & Bassa \\
\bottomrule
\end{longtable}

\subsection{Resource Allocation Plan}
\begin{itemize}
  \item VM-DC-01: 2 vCPU, 4 GB RAM (dynamic), disco OS 60 GB.
  \item VM-LINUX-01: 4 vCPU, 8 GB RAM (dynamic), disco OS 50 GB.
  \item VM-SQL-01: 4 vCPU, 8 GB RAM (static), disco 100 GB.
\end{itemize}

\subsection{Services Catalog}
\subsubsection*{VM-LINUX-01 (Docker Environment)}
\begin{itemize}
  \item Plex (TCP 32400): streaming video.
  \item Navidrome (TCP 4533): streaming audio self-hosted.
  \item Immich (TCP 2283): backup foto/video mobile.
  \item Cloudflare Tunnel (\texttt{cloudflared}): esposizione sicura servizi web.
\end{itemize}

\subsubsection*{VM-SQL-01 (Development Lab)}
\begin{itemize}
  \item SQL Server (TCP 1433), autenticazione mista per scenari di test.
  \item IIS (HTTP/HTTPS) per portfolio e applicazioni interne.
\end{itemize}

\subsection{Automatic Virtual Machine Activation (AVMA)}
Le VM Windows possono essere attivate via AVMA utilizzando la licenza Datacenter dell'host fisico.

\section{Network \& Connectivity}

\subsection{IP Addressing Plan (IPv4)}
\begin{longtable}{>{\raggedright\arraybackslash}p{0.24\textwidth} >{\raggedright\arraybackslash}p{0.15\textwidth} >{\raggedright\arraybackslash}p{0.24\textwidth} >{\raggedright\arraybackslash}p{0.26\textwidth}}
\toprule
\textbf{Nodo} & \textbf{IP statico} & \textbf{Ruolo} & \textbf{Servizi principali} \\
\midrule
\endfirsthead
\toprule
\textbf{Nodo} & \textbf{IP statico} & \textbf{Ruolo} & \textbf{Servizi principali} \\
\midrule
\endhead
\texttt{HV-NODE-01} & 192.168.1.10 & Hypervisor/Storage & RDP, SMB \\
\texttt{DC-01} & 192.168.1.11 & AD DS, DNS, DHCP & DNS, LDAP \\
\texttt{LINUX-01} & 192.168.1.12 & Docker Host & Plex, Immich, Navidrome \\
\texttt{SQL-01} & 192.168.1.13 & DB/Web & SQL Server, IIS \\
\bottomrule
\end{longtable}

\subsection{DNS Strategy}
\begin{itemize}
  \item DNS autoritativo interno: \texttt{VM-DC-01}.
  \item Forwarders: resolver pubblici affidabili per richieste esterne.
  \item Split DNS per servizi con endpoint interni/esterni distinti.
\end{itemize}

\subsection{Remote Access (Private Plane)}
\begin{itemize}
  \item Tecnologia: Tailscale (mesh VPN WireGuard).
  \item Nodo gateway: \texttt{VM-LINUX-01} come subnet router.
  \item Use case: RDP amministrativo, SMB remoto protetto, accesso servizi LAN.
\end{itemize}

\subsection{Public Exposure (Application Plane)}
\begin{itemize}
  \item Tecnologia: Cloudflare Tunnel per pubblicazione L7.
  \item Nessun port-forwarding diretto sul router ISP.
  \item Mapping pubblico selettivo verso servizi interni necessari.
\end{itemize}

\subsection{Firewall Rules}
\begin{itemize}
  \item Consentire solo porte necessarie e solo da subnet/range previsti.
  \item Separare policy interne (LAN/VPN) e policy pubbliche (Cloudflare).
  \item Applicare principio \textit{least privilege} su host e guest OS.
\end{itemize}

\section{Backup \& Disaster Recovery (DR)}

\subsection{Backup Strategy (3-2-1)}
\begin{itemize}
  \item 3 copie: produzione, backup locale, backup off-site.
  \item 2 media: storage interno + USB/cloud.
  \item 1 copia off-site per dati critici (foto/documenti/codice).
\end{itemize}

Tooling principale:
\begin{itemize}
  \item Veeam Backup \& Replication Community Edition.
  \item Volume Shadow Copy (VSS) per recupero rapido file.
\end{itemize}

\subsection{Backup Jobs Configuration}
\begin{longtable}{>{\raggedright\arraybackslash}p{0.22\textwidth} >{\raggedright\arraybackslash}p{0.24\textwidth} >{\raggedright\arraybackslash}p{0.18\textwidth} >{\raggedright\arraybackslash}p{0.12\textwidth} >{\raggedright\arraybackslash}p{0.16\textwidth}}
\toprule
\textbf{Job} & \textbf{Sorgente} & \textbf{Destinazione} & \textbf{Schedulazione} & \textbf{Retention} \\
\midrule
\endfirsthead
\toprule
\textbf{Job} & \textbf{Sorgente} & \textbf{Destinazione} & \textbf{Schedulazione} & \textbf{Retention} \\
\midrule
\endhead
BKP-VMs-Daily & VM DC/LINUX/SQL & Repository USB & Giornaliero 02:00 & 14 giorni \\
BKP-Files-Crit & \texttt{D:\textbackslash{}FAMILY\_BACKUP}, \texttt{D:\textbackslash{}DEV}, \texttt{D:\textbackslash{}APP\_DATA} & USB + sync cloud & Giornaliero 04:00 & 30 giorni \\
BKP-Media-Wkly & \texttt{D:\textbackslash{}MEDIA} & USB & Settimanale & 2 versioni \\
VSS-Snapshots & Volume \texttt{D:} & Shadow storage locale & Ogni 6 ore & Max 10\% \\
\bottomrule
\end{longtable}

\subsection{Runbooks di Emergenza}
\paragraph{Scenario A: Guasto HDD}
Identificazione disco degradato, sostituzione fisica, reintroduzione nel pool e avvio processo di repair del virtual disk.

\paragraph{Scenario B: Corruzione OS Host}
Reinstallazione host su NVMe, riaggancio Storage Spaces, import VM da \texttt{D:\textbackslash{}DEV\textbackslash{}VMs}, riconfigurazione rete/share.

\paragraph{Scenario C: Ransomware}
Isolamento immediato rete, tentativo restore da VSS, fallback restore da Veeam, bonifica host e ripristino controllato.

\section{Conclusions \& Critical Takeaways}

\subsection{Executive Summary}
HP Fortress dimostra che una piattaforma enterprise-like può essere realizzata con budget contenuto, mantenendo alta disponibilità pratica, segregazione dei ruoli e governance dei dati.

\subsection{Golden Rules}
\begin{enumerate}
  \item RAID non è backup.
  \item Le licenze e le chiavi vanno custodite fuori dal server.
  \item Host fisico dedicato a hypervisor/storage, applicazioni nelle VM.
  \item Accesso remoto solo tramite canali sicuri (Tailscale/Cloudflare).
\end{enumerate}

\subsection{Maintenance Schedule}
\begin{itemize}
  \item \textbf{Mensile}: patch OS host/guest, update container, verifica stato dischi.
  \item \textbf{Trimestrale}: test di restore backup e pulizia fisica del nodo.
  \item \textbf{Annuale}: review licenze/certificati e aggiornamento documentazione.
\end{itemize}

\subsection{Future Roadmap}
\begin{itemize}
  \item Upgrade networking multi-gig.
  \item Espansione RAM per laboratori più pesanti (es. Kubernetes/EVE-NG).
  \item Maggiore automazione monitoring e alerting.
\end{itemize}

\section{Secure Remote Access \& Zero Trust Architecture}

\subsection{Philosophy: Invisible Fortress}
Approccio \textit{zero open ports}: nessun ingresso diretto dal router; solo tunnel cifrati pre-autenticati.

\subsection{Split-Tunnel Topology}
\begin{itemize}
  \item \textbf{Management Plane}: Tailscale (accesso L3 completo alla LAN per amministrazione).
  \item \textbf{Application Plane}: Cloudflare Tunnel (pubblicazione selettiva L7 di servizi web).
\end{itemize}

\subsection{Channel A: Tailscale}
\begin{itemize}
  \item Nodo subnet-router: \texttt{VM-LINUX-01}.
  \item Use case: RDP, SMB, SQL management da remoto in tunnel WireGuard.
  \item ACL opzionali per ridurre superficie di accesso dei device non amministrativi.
\end{itemize}

\subsection{Channel B: Cloudflare Tunnel}
\begin{itemize}
  \item Tier pubblico (es. portfolio) con WAF e HTTPS always-on.
  \item Tier privato (es. Immich/Portainer/Stirling PDF) protetto da Cloudflare Access.
  \item Policy differenziate per sottodominio e rischio applicativo.
\end{itemize}

\subsection{Security Hardening Checklist}
\begin{itemize}
  \item WAF geofencing su aree non pertinenti.
  \item HSTS abilitato sui domini pubblici.
  \item Verifica scadenza chiavi Tailscale sul nodo server.
  \item Fallback firewall locale con allowlist minimali.
\end{itemize}

\subsection{Access DR}
In caso di fault del gateway Linux (Tailscale/Cloudflare), prevedere accesso fisico locale come canale di emergenza e procedura documentata di recovery dell'accesso remoto.

\end{document}
